\section{Discussion and Limitations}\label{sec:attack:discussion}

%% BIG CHANGE: polish; restructure the argument and pacing, retaining the stated scope.
Unpaired matching improves reconstruction word-F1 when the affine map is fitted on a small leak.
With \vf{90} leaked rows, mean word-F1 on \vf{4{,}500} held-out prompts rises from \vf{0.354} to \vf{0.766} on Corpus A and from \vf{0.102} to \vf{0.321} on legal128w, averaged over three seeds. % src: results/v7_2/main/w46-confirmation-s{1,2,3}.json; results/legal128w/v7_2/main/w46-confirmation-s{1,2,3}.json
Both gains meet H2, including its required interval lower bound of \vf{0.05}.
The ablations delimit this result: the gain shrinks as the leak grows and depends on the temperature, while the tested blind maps exceed their controls by at most \vf{0.025} word-F1.

The exploratory rewriting experiment yields a different comparison.
A second-round language model exceeds the mean of the three decoders by \vf{0.12} to \vf{0.26} word-F1 in three conditions, but exceeds the best decoder by only \vf{0.01} to \vf{0.02}, without an interval for that difference. % src: results/*/v7_2/main/framework/attacker/*-haiku.json
This result rests on one seed and selection prompts; it does not establish a reliable improvement over the best decoder.

\paragraph{Resources and alignment.} The attack requires the frozen hub and its decoders, the hidden states of the target model for prompts without text, the attacker's own text of the same domain, and a leak of rows.
%% BIG CHANGE: polish; reorganize explanation without changing the reported evidence.
The leak supplies the paired correspondences for fitting the map; constructing those rows requires the target tokenizer for the leaked prompts.
The experiments also select settings on \vf{500} prompts with known text (Section~\ref{sec:attack:threat}).
The unpaired pools carry no pairing between hidden states and text, yet they account for the difference between the supervised and the hybrid fit.
%% BIG CHANGE: polish; reword or relocate this passage for precision and flow.
The second round receives only the decoders' reconstructions and the kind of text.

%% BIG CHANGE: polish; reword or relocate this passage for precision and flow.
\paragraph{What the reconstructions contain.} Word-F1 measures overlap between the reconstructed and reference words.
On legal text the reconstructions keep the frame of a passage and substitute names and numbers (Table~\ref{tab:attack:legal-examples}), and after the second round entity recall stays at \vf{0.37} for a budget of \vf{10} (Section~\ref{sec:attack:rewriting}).
On legal128w the hybrid fit also has the higher teacher-forced cross-entropy (Section~\ref{sec:attack:teacher}).
%% BIG CHANGE: polish; reword or relocate this passage for precision and flow.
Exact match under the hybrid fit is \vf{0.21} on Corpus A and rounds to zero on legal128w (Section~\ref{sec:attack:hybrid}).
The results therefore establish a gain in word overlap; they do not establish exact recovery, recovery of the facts of a prompt, or an amount of leaked information.

\paragraph{Limitations.} Six limitations bound the claims.
\begin{itemize}
%% BIG CHANGE: polish; reword or relocate this passage for precision and flow.
\item \textbf{Held-out prompts.} Only the main comparison and two ablations evaluate the \vf{4{,}500} held-out prompts.
Most other scores come from the \vf{500} selection prompts, which also selected the settings; the second round evaluates \vf{450} of them, and one blind comparison evaluates all \vf{5{,}000} validation prompts.
The held-out prompts are held out from the attack settings, not from the hub checkpoint, which all validation prompts selected.
%% BIG CHANGE: polish; reword or relocate this passage for precision and flow.
The test splits have not been evaluated.
\item \textbf{Seeds.} The main comparison and the Corpus A budget curve rest on three seeds; the second round and the legal128w teacher-forced scores rest on one, and each ablation table states its seeds.
The intervals describe variation over prompts, not over seeds.
\item \textbf{Several settings.} The experiments differ in supervision fraction, temperature, batch size, and form of the leak (Table~\ref{tab:attack:settings}), so their scores do not combine into one curve.
\item \textbf{Settings chosen with known text.} The supervision fraction and the temperature were selected by word-F1 on prompts whose text is known.
An attacker without such prompts would have to fix them in another way, and no run measures the cost.
\item \textbf{Same-domain pools.} Both unpaired pools and the leak come from the training prompts of the corpus under attack.
The results do not show how the attack behaves when the attacker's text comes from another domain, or when the unpaired pool is small.
\item \textbf{Models and layer.} The results cover one block of two target models of about one billion parameters, three hub models, and one second-round language model.
They do not cover a deployed system, a larger target model, a quantized or noised hidden state, or a defense.
\end{itemize}

The leak supplies correspondences that constrain the affine map, and the matching term compares unpaired representations; the experiments do not establish that either source determines the map alone.
The experiments quantify reconstruction under the stated attack resources; an information-theoretic account remains future work (Chapter~\ref{ch:future}).
